% Appendix A: derivation of the gap rule (Proposition 1), with corollaries and
% examples. Source: for-claude/gap-rule-proof-v1.md (2026-08-14), converted 2026-09-08.
% Self-contained: states every assumption used; no KKT machinery (Lemma 3 does that
% work via the transfer argument and the concavity gradient inequality). Numerical
% verification: for-claude/verify_gap_rule.py (81 trials vs scipy SLSQP, machine
% precision; see the md's note 5).
\section{Derivation of the gap rule}\label{app:gap-rule}

\subsection*{The within-round allocation problem}

A population of $n$ researchers. Researcher $i$ has capability $K_i$ and baseline resources $R_{i0}$, both fixed within the round. A grant $g_i \geq 0$ adds to researcher $i$'s resources for the round, so that their expected output is $A\,\Lambda(K_i,\, R_{i0} + g_i)$ with $\Lambda(K, R) = 2KR/(K + R)$, the production function of \cref{sec:model}, and $A > 0$ the common productivity scale. Write
\[
\varphi_i(g) = \frac{2A\,K_i\,(R_{i0} + g)}{K_i + R_{i0} + g}
\]
for researcher $i$'s expected output as a function of their grant alone.

Throughout, $A > 0$, $K_i > 0$ and $R_{i0} \geq 0$ for every $i$, and the round's budget is $B > 0$. Nothing else about the model is used, so the result holds for every realized population and every round's budget; across rounds, the constant $c$ and the funded set are recomputed from the new population state.

The funder observes $(K_i, R_{i0})_{i=1}^n$ and divides the budget so as to maximize expected output:
\[
(\mathrm{P}) \qquad \max_{g_1, \ldots, g_n \geq 0} \; \sum_{i=1}^n \varphi_i(g_i) \quad \text{subject to} \quad \sum_{i=1}^n g_i \leq B.
\]

\subsection*{Three lemmas}

\begin{lemma}[the return to a grant]\label{lem:return}
For each $i$, on $g \in [0, \infty)$: (i) $\varphi_i$ is strictly increasing, with $\varphi_i'(g) = 2A\,K_i^2/(K_i + R_{i0} + g)^2 > 0$; (ii) $\varphi_i$ is strictly concave: $\varphi_i''(g) = -\,4A\,K_i^2/(K_i + R_{i0} + g)^3 < 0$; (iii) $\varphi_i'(g) < 2A$ whenever $R_{i0} + g > 0$, and $\varphi_i(g) \uparrow 2A K_i$ as $g \to \infty$.
\end{lemma}

\begin{proof}
Write $S = K_i + R_{i0} + g$. Then $\varphi_i(g) = 2AK_i(S - K_i)/S = 2AK_i - 2AK_i^2/S$, and (i) and (ii) follow by differentiating in $g$ (note $dS/dg = 1$). For (iii): $\varphi_i'(g) = 2A\,K_i^2/S^2 < 2A$ exactly when $S > K_i$, that is, when $R_{i0} + g > 0$; and $\varphi_i(g) = 2AK_i - 2AK_i^2/S \to 2AK_i$ as $S \to \infty$.
\end{proof}

\Cref{lem:return} says that the marginal value of a dollar to researcher $i$ increases with their capability and decreases with the resources they already hold, that each dollar granted lowers the value of the next, and that expected output never exceeds $2AK_i$.

\begin{lemma}[existence, uniqueness, exhaustion]\label{lem:existence}
Problem $(\mathrm{P})$ has exactly one solution $g^* = (g_1^*, \ldots, g_n^*)$, and it exhausts the budget: $\sum_i g_i^* = B$.
\end{lemma}
\begin{proof}
The feasible set is compact and convex and the objective $f(g) = \sum_i \varphi_i(g_i)$ is continuous and, by \cref{lem:return}(ii), strictly concave, so a maximizer exists and is unique. If $\sum_i g_i^* < B$, increasing any one $g_i^*$ slightly stays feasible and raises $f$ (\cref{lem:return}(i)), contradicting optimality.
\end{proof}
\begin{lemma}[equal marginal value]\label{lem:emv}
A feasible allocation $g$ with $\sum_i g_i = B$ solves $(\mathrm{P})$ if and only if there is a $\nu > 0$ such that
\[
\varphi_i'(g_i) = \nu \;\text{ for every } i \text{ with } g_i > 0, \qquad \varphi_i'(0) \leq \nu \;\text{ for every } i \text{ with } g_i = 0.
\]
\end{lemma}

\begin{proof}
\emph{Necessity.} Let $g^*$ solve $(\mathrm{P})$. Since $B > 0$ is exhausted (\cref{lem:existence}), some researcher is funded: $g_i^* > 0$ for some $i$. Fix any such $i$ and any $j \neq i$, and for $0 < \delta \leq g_i^*$ consider the transfer that replaces $g_i^*$ by $g_i^* - \delta$ and $g_j^*$ by $g_j^* + \delta$, leaving the rest unchanged. The result is feasible, so optimality gives
\[
\varphi_j(g_j^* + \delta) - \varphi_j(g_j^*) \;\leq\; \varphi_i(g_i^*) - \varphi_i(g_i^* - \delta).
\]
Dividing by $\delta$ and letting $\delta \downarrow 0$ yields $\varphi_j'(g_j^*) \leq \varphi_i'(g_i^*)$. If $j$ is also funded, the same argument with $i$ and $j$ exchanged gives the reverse inequality, so all funded researchers share a common marginal value; call it $\nu = \varphi_i'(g_i^*) > 0$ (positivity by \cref{lem:return}(i)). For unfunded $j$, the displayed inequality reads $\varphi_j'(0) \leq \nu$.

\emph{Sufficiency.} Let $g$ satisfy the condition with multiplier $\nu$, and let $h$ be any feasible allocation. Concavity of $\varphi_i$ gives, for every $i$, $\varphi_i(h_i) \leq \varphi_i(g_i) + \varphi_i'(g_i)\,(h_i - g_i)$. For funded $i$, $\varphi_i'(g_i) = \nu$, so $\varphi_i'(g_i)(h_i - g_i) = \nu\,(h_i - g_i)$. For unfunded $i$, $g_i = 0$ and $h_i - g_i = h_i \geq 0$, so $\varphi_i'(0)\,(h_i - g_i) \leq \nu\,(h_i - g_i)$. Summing over $i$:
\[
f(h) \;\leq\; f(g) + \nu \sum_i (h_i - g_i) \;=\; f(g) + \nu\Big(\sum_i h_i - B\Big) \;\leq\; f(g),
\]
since $h$ is feasible and $\nu > 0$. So $g$ solves $(\mathrm{P})$.
\end{proof}

\Cref{lem:emv} is the formal content of the sentence in \cref{sec:optimal}: an allocation is optimal exactly when no dollar can be moved to a researcher who can produce more value with it.

\subsection*{The proposition}

\begin{proprestate}[the gap rule]
The unique solution of $(\mathrm{P})$ is
\[
g_i^* = \max(c\,K_i - R_{i0},\, 0), \qquad i = 1, \ldots, n,
\]
where $c > 0$ is the unique constant satisfying the budget identity $\sum_{i=1}^n \max(c\,K_i - R_{i0},\, 0) = B$. Equivalently, $c = \sqrt{2A/\nu} - 1$, where $\nu \in (0, 2A)$ is the common marginal value of \cref{lem:emv}.
\end{proprestate}

\begin{proof}
Let $g^*$ be the unique solution (\cref{lem:existence}) and $\nu$ the multiplier that \cref{lem:emv} associates with it. The proof proceeds in four steps.

\emph{Step 1: $\nu \in (0, 2A)$, so $c = \sqrt{2A/\nu} - 1$ is well defined and positive.} Positivity of $\nu$ is part of \cref{lem:emv}. For the upper bound: some researcher $i$ is funded, and for them $R_{i0} + g_i^* \geq g_i^* > 0$, so \cref{lem:return}(iii) gives $\nu = \varphi_i'(g_i^*) < 2A$. Then $2A/\nu > 1$, so $c = \sqrt{2A/\nu} - 1 > 0$.

\emph{Step 2: funded researchers receive exactly their gap.} Let $g_i^* > 0$. By \cref{lem:emv}, $2A\,K_i^2/(K_i + R_{i0} + g_i^*)^2 = \nu$. Solving, and taking the positive square root (both sides of the next display are positive),
\[
K_i + R_{i0} + g_i^* = K_i\,\sqrt{2A/\nu} = (1 + c)\,K_i,
\]
so $g_i^* = c\,K_i - R_{i0}$. Since $g_i^* > 0$, this equals $\max(c\,K_i - R_{i0},\, 0)$.

\emph{Step 3: unfunded researchers have no gap.} Let $g_i^* = 0$. By \cref{lem:emv}, $2A\,K_i^2/(K_i + R_{i0})^2 = \varphi_i'(0) \leq \nu$, so $K_i + R_{i0} \geq K_i\sqrt{2A/\nu} = (1 + c)\,K_i$, that is, $R_{i0} \geq c\,K_i$. Then $\max(c\,K_i - R_{i0},\, 0) = 0 = g_i^*$.

Steps 2 and 3 establish $g_i^* = \max(c\,K_i - R_{i0}, 0)$ for every $i$, and the budget identity is exhaustion (\cref{lem:existence}). It remains to show $c$ is the only constant satisfying that identity.

\emph{Step 4: the budget identity pins $c$ uniquely.} Define $G(c) = \sum_{i=1}^n \max(c\,K_i - R_{i0},\, 0)$ for $c \geq 0$. $G$ is continuous and nondecreasing, with $G(0) = 0$ and $G(c) \to \infty$ as $c \to \infty$. Moreover $G$ is strictly increasing wherever it is positive: if $G(c) > 0$, then $c\,K_j > R_{j0}$ for some $j$, and for any $c' > c$, $G(c') \geq G(c) + K_j\,(c' - c) > G(c)$. Now suppose $c_1 < c_2$ both satisfy $G(c) = B$. Since $B > 0$, $G(c_1) > 0$, so $G(c_2) > G(c_1) = B$, a contradiction.
\end{proof}

\subsection*{Consequences}

\begin{corollary}[funding frontier and budget comparative statics]\label{cor:frontier}
(i) Researcher $i$ is funded, $g_i^* > 0$, exactly when $K_i > R_{i0}/c$. (ii) As a function of the budget $B$, the constant $c(B)$ is continuous and strictly increasing; hence the funded set is nondecreasing in $B$, every funded researcher's grant $c(B)\,K_i - R_{i0}$ is strictly increasing in $B$, and the common marginal value $\nu(B) = 2A/(1 + c(B))^2$ is strictly decreasing in $B$. Both parts follow from the formula and Step 4.
\end{corollary}
\begin{corollary}[the value of targeting vanishes in the budget]\label{cor:targeting-vanishes}
Let $f^*(B)$ denote the optimal value of $(\mathrm{P})$ and $f^u(B) = \sum_i \varphi_i(B/n)$ the value of uniform funding. Then $0 \leq f^*(B) - f^u(B)$ for every $B$, and $f^*(B) - f^u(B) \to 0$ as $B \to \infty$.
\end{corollary}

\begin{proof}
The first inequality holds because uniform funding is feasible. For the second: by \cref{lem:return}(iii), $\varphi_i(g) < 2AK_i$ for every $g$, so $f^*(B) \leq 2A\sum_i K_i$; and $f^u(B) = \sum_i \varphi_i(B/n) \to 2A\sum_i K_i$ as $B \to \infty$. Hence $0 \leq f^*(B) - f^u(B) \leq 2A\sum_i K_i - f^u(B) \to 0$.
\end{proof}

Since uniform funding's own gain over no funding does not vanish (it approaches $2A\sum_i K_i - \sum_i \varphi_i(0) > 0$), the same conclusion holds for the value of targeting measured relative to that gain, the quantity plotted in Figure~\ref{fig:targeting-value}.

\begin{example}[funding the track record can produce less than uniform funding]\label{ex:track-record}
Two researchers, $A = 1$, budget $B = 2$. Researcher 1: $K_1 = 8$, $R_{10} = 8$ (productive and well resourced; expected output $\varphi_1(0) = 8$). Researcher 2: $K_2 = 4$, $R_{20} = 1$ (capable but resource-starved; $\varphi_2(0) = 8/5$). Allocating in proportion to expected output gives $g = (5/3,\, 1/3)$ and total output $570/53 \approx 10.75$. Uniform funding, $g = (1, 1)$, gives $568/51 \approx 11.14$. The gap rule ($c = 3/4$) gives $g^* = (0, 2)$ and $80/7 \approx 11.43$.
\end{example}

\begin{example}[funding the under-resourced can produce less than uniform funding]\label{ex:scarcity}
Two researchers, $A = 1$, budget $B = 2$. Researcher 1: $K_1 = 10$, $R_{10} = 5$. Researcher 2: $K_2 = 1/2$, $R_{20} = 0$ (the scarcest resources in the population). Directing the budget to the most resource-starved researcher gives total output $112/15 \approx 7.47$. Uniform funding gives $49/6 \approx 8.17$. The gap rule ($c = 2/3$) gives $g^* = (5/3,\, 1/3)$ and $42/5 = 8.4$.
\end{example}
